% !TEX root = ../main.tex
% !TEX program = lualatex
\documentclass[../main.tex]{subfiles}
\IfSubfilesClassLoaded{\externaldocument{\subfix{../build/main}}}{}
% =============================================================================
% File: 27G_CPI_260.tex
% =============================================================================
\begin{document}
\IfSubfilesClassLoaded{\chapter{California Psychological Inventory 260}}{}
\section{California Psychological Inventory 260}

\subsection{Summary}
\begin{description}
  \item[\textbf{Purpose.}] CPI 260 provides objective, comparative feedback on leadership tendencies and comfort zones to highlight strengths and potential blind spots~\autocite{edo-5-1-9-cpi-260-2025}.
  \item[\textbf{Reports.}] The Client Feedback Report shows scores across scales and the Coaching Report for Leaders translates those scores into leadership development insights~\autocite{edo-5-1-9-cpi-260-2025}.
  \item[\textbf{Action focus.}] Results are hypotheses, not truths; leaders are expected to interpret, corroborate, and build an action plan~\autocite{edo-5-1-9-cpi-260-2025}.
\end{description}

\subsection{History and Methodology}
\begin{description}
  \item[\textbf{Origins.}] The CPI family dates to 1956 (CPI 480); CPI 260 launched in 2002 and has been used in the senior course since 2003~\autocite{edo-5-1-9-cpi-260-2025}.
  \item[\textbf{Purpose.}] CPI 260 describes you from the perspective of knowledgeable, objective others, capturing comfort zones that may reveal blind spots~\autocite{edo-5-1-9-cpi-260-2025}.
  \item[\textbf{Data foundation.}] Scale definitions rely on observed trends and comparative analysis (not theory alone), with demographic screening to retain the most discriminating scales~\autocite{edo-5-1-9-cpi-260-2025}.
\end{description}

\subsection{Client Feedback Report}
\begin{description}
  \item[\textbf{Metrics baseline.}] The Client Feedback Report displays CPI 260 scores across 29 scales and compares them to the U.S. adult population~\autocite{edo-5-1-9-cpi-260-2025}.
  \item[\textbf{Screening results.}] Response screening flags overly favorable, overly negative, or random patterns that require further investigation~\autocite{edo-5-1-9-cpi-260-2025}.
\end{description}

\subsection{Lifestyle Scales and Management Styles}
\begin{description}
  \item[\textbf{Orientation.}] Lifestyle scales capture how you engage others and relate to norms, ranging from interactive and expressive to reflective and reserved~\autocite{edo-5-1-9-cpi-260-2025}.
  \item[\textbf{Quadrants.}] Management styles map to implementer, supporter, innovator, and visualizer quadrants with distinct strengths and risks in change and execution~\autocite{edo-5-1-9-cpi-260-2025}.
  \item[\textbf{Satisfaction scale.}] Satisfaction strongly influences overall results; higher scores tend to raise other scales while lower scores suppress them~\autocite{edo-5-1-9-cpi-260-2025}.
\end{description}

\subsection{Scale Families}
\begin{description}
  \item[\textbf{Dealing with others.}] Dominance, sociability, social presence, self-acceptance, independence, and empathy compare you to a national norm~\autocite{edo-5-1-9-cpi-260-2025}.
  \item[\textbf{Self-management.}] Responsibility, social conformity, self-control, good impression, communality, well-being, and tolerance track self-regulation tendencies~\autocite{edo-5-1-9-cpi-260-2025}.
  \item[\textbf{Motivations and thinking.}] Achievement via conformance, achievement via independence, and conceptual fluency provide insight into how you pursue goals~\autocite{edo-5-1-9-cpi-260-2025}.
  \item[\textbf{Personal characteristics.}] Insightfulness, flexibility, and sensitivity signal how you process interpersonal and organizational dynamics~\autocite{edo-5-1-9-cpi-260-2025}.
  \item[\textbf{Work-related measures.}] Managerial potential, work orientation, creative temperament, leadership, amicability, and law enforcement orientation capture applied leadership behaviors~\autocite{edo-5-1-9-cpi-260-2025}.
\end{description}

\subsection{Interpreting Results}
\begin{description}
  \item[\textbf{Let\'s talk approach.}] Treat results as approximations, compare to comfort zones, corroborate with other feedback, and define follow-through actions~\autocite{edo-5-1-9-cpi-260-2025}.
  \item[\textbf{No ideal profile.}] There is no perfect set of results; each situation is unique and change is possible and worthwhile~\autocite{edo-5-1-9-cpi-260-2025}.
  \item[\textbf{Johari window.}] Feedback reduces blind spots and improves self-awareness in leadership behavior~\autocite{edo-5-1-9-cpi-260-2025}.
\end{description}

\subsection{Coaching Report for Leaders}
\begin{description}
  \item[\textbf{Purpose.}] The Coaching Report for Leaders converts CPI scale data into leadership strengths, developmental opportunities, and areas to evaluate further~\autocite{edo-5-1-9-cpi-260-2025}.
  \item[\textbf{Executive norm group.}] Leadership characteristics compare CPI scale pairs to an executive norm group to flag strengths, developmental opportunities, or undetermined areas~\autocite{edo-5-1-9-cpi-260-2025}.
  \item[\textbf{Next steps.}] Focus on two or three key priorities, leverage strengths, and address blind spots with targeted action~\autocite{edo-5-1-9-cpi-260-2025}.
\end{description}

\subsection{Quick Review}
\begin{itemize}
  \item CPI 260 is a self-awareness tool that highlights leadership tendencies and blind spots~\autocite{edo-5-1-9-cpi-260-2025}.
  \item The Client Feedback Report provides scale scores; the Coaching Report for Leaders turns them into development guidance~\autocite{edo-5-1-9-cpi-260-2025}.
  \item Results are hypotheses that require corroboration and action planning~\autocite{edo-5-1-9-cpi-260-2025}.
\end{itemize}

\IfSubfilesClassLoaded{
  \RenewDocumentCommand{\entryname}{}{\textbf{\color{Modern} Acronym}}
  \RenewDocumentCommand{\descriptionname}{}{\textbf{\color{Modern} Definition}}
  \printnoidxglossary[
    type=\acronymtype,
    title=Acronyms,
    style=long-booktabs
  ]}{}
\end{document}
